\chapter[Составной тетраэдрический порядок]{Составной тетраэдрический
порядок и решение калибровочно-кадровой развилки}

Предыдущая глава оставила две формально допустимые, но физически различные
реализации тетраэдрической фазы. Независимое поле спина три сохраняет
остаточную калибровочную группу $A_4$, но ещё не принадлежит конечному
родителю. Четырёхосный кадр уже построен из аффинного носителя, но его
левый калибровочный стабилизатор тривиален. Настоящий гейт должен решить,
можно ли получить тетраэдрический порядок как составной момент, не
подменяя глобальную или граничную симметрию калибровочной.

\section{Литературная граница составного порядка}

В теории тетраэдрических жидкокристаллических фаз известна возможность
состояния, в котором векторный и квадрупольный параметры порядка равны
нулю, а полностью симметричный бесследовый тензор третьего ранга ненулевой;
см. \path{arXiv:cond-mat/0008325}, \path{arXiv:cond-mat/0205171} и
\path{arXiv:2210.00575}. Поэтому ненулевой третий момент без
конденсации ориентированного кадра является признанным физическим
механизмом, а не только формальной операцией над тензорами.

Литература, однако, не решает родительскую задачу проекта. Она допускает
тензор третьего ранга как эффективный параметр порядка, но не доказывает,
что он возникает как поле или кривизна данного конечного оператора.

\section{Точный тетраэдрический третий момент}

Пусть $n_a$, $a=1,\ldots,4$, --- четыре единичные вершины правильного
тетраэдра и $p_a=1/4$. Определим моменты
\begin{equation}
 m_i=\sum_ap_a n_{a,i},\qquad
 C_{ij}=\sum_ap_a n_{a,i}n_{a,j},\qquad
 \Theta_{ijk}=\sum_ap_a n_{a,i}n_{a,j}n_{a,k}.
 \label{eq:v6-tetrahedral-moment-hierarchy}
\end{equation}
Прямая проверка даёт
\begin{equation}
 m=0,\qquad C=\frac13I_3,\qquad
 \|\Theta\|^2=\frac29.
 \label{eq:v6-pure-third-rank-order}
\end{equation}
Следовательно, среда не выбирает ни направленный вектор, ни нематическую
ось, но уже различает тетраэдрическую ориентацию третьим моментом.

Двенадцать собственных вращений тетраэдра сохраняют $\Theta$. Ранг её
непрерывной $SO(3)$-орбиты равен трём, а массовая матрица одной семейной
связности имеет спектр
\begin{equation}
 \Spec M_\Theta=\left\{\frac89,\frac89,\frac89\right\}.
 \label{eq:v6-composite-third-moment-gauge-mass}
\end{equation}
Таким образом, составной третий момент геометрически даёт именно
$SO(3)/A_4$ и положительную массу всех трёх непрерывных направлений.

\begin{proposition}[Тетраэдрический составной момент]
Равномерное состояние на четырёх тетраэдрических вершинах имеет нулевые
первый и бесследовый второй моменты, но ненулевой третий момент со
стабилизатором $A_4$. Поэтому независимый линейный конденсат кадра не
является математически обязательным для тетраэдрического порядка.
\end{proposition}

Это снимает только геометрическую часть развилки. Осталось проверить,
является ли $\Theta$ динамическим составным полем текущего родителя.

\section{Представительный состав конечного родителя}

Семейный триплет есть представление $V_1$ спина один. Его операторный
алгебраический носитель раскладывается как
\begin{equation}
 \operatorname{End}(V_1)=V_0\oplus V_1\oplus V_2.
 \label{eq:v6-end-triplet-spin-bound}
\end{equation}
Именно поэтому матричные семейные поля дают скаляр, связность и
пятикомпонентный параметр $Q$, но не семикомпонентное поле $V_3$.

Аффинный четырёхточечный носитель имеет
$V_4=V_0\oplus V_1$. Поэтому стрелка кадра удовлетворяет
\begin{equation}
 \operatorname{Hom}(V_4,V_1)
 =V_0\oplus2V_1\oplus V_2,
 \label{eq:v6-frame-arrow-decomposition}
\end{equation}
и также не содержит спина три.

Для полного семейного модуля
\begin{equation}
 K_{\rm fam}=V_4\oplus V_1\oplus V_1
\end{equation}
машинный разбор даёт
\begin{equation}
 \operatorname{End}(K_{\rm fam})
 =10V_0\oplus15V_1\oplus9V_2,
 \qquad \dim\operatorname{End}(K_{\rm fam})=100.
 \label{eq:v6-full-parent-spin-ledger}
\end{equation}
Кратность $V_3$ точно равна нулю. Умножение операторов не помогает:
обычные матричные многочлены остаются в той же конечной алгебре
$\operatorname{End}(K_{\rm fam})$.

Напротив, внешний симметрический куб триплета содержит
\begin{equation}
 \operatorname{Sym}^3(V_1)=V_3\oplus V_1.
 \label{eq:v6-external-symmetric-cube}
\end{equation}
Искомая семёрка появляется здесь ровно один раз. Но внешний тензорный куб
является новым носителем относительно текущей операторной алгебры. Его
кинетическая норма и кривизна не следуют автоматически из следа уже
существующего блока.

\begin{theorem}[Операторный запрет текущего родителя]
В $\operatorname{End}(K_{\rm fam})$ и во всех его обычных матричных многочленах нет
представления спина три. Поэтому тетраэдрический третий момент существует
как точная геометрическая и эффективная переменная, но ещё не является
динамическим полем данного конечного родителя. Для него необходим внешний
носитель $\operatorname{Sym}^3(V_1)$ или эквивалентный новый узел конечной геометрии.
\end{theorem}

\section{Почему упорядоченный кадр не заменяет этот носитель}

Строгий аффинный слой ранее установил, что полная $S_4$-симметрия меню
является условным семейным постулатом, а не выведенной локальной
калибровочной группой. Кроме того, для коммутативного четырёхточечного
узла перестановки являются внешними автоморфизмами, а не внутренними
калибровочными преобразованиями.

Поэтому равенство
\begin{equation}
 R_\sigma X_\star=X_\star U_\sigma
\end{equation}
сохраняет диагональную глобальную или гранично фиксированную $A_4$, но не
остаточную подгруппу одного левого калибровочного $SO(3)$. Сам полный кадр
$X_\star$ имеет тривиальный левый стабилизатор и полностью хиггсует эту
связность.

Это стандартная граница калибровочно-семейного сцепления: после
конденсации бифундаментального поля калибровочная группа может быть
полностью нарушена, тогда как диагональная группа остаётся глобальной.
Она становится диагональной калибровочной группой только при калибровке
обоих множителей, то есть после введения второй связности.

\begin{nogocriterion}[Запрет скрытой калибровки перестановок]
Нельзя считать правое действие $S_4$ локальной калибровочной
эквивалентностью только потому, что оно переставляет четыре примитивные
точки. Для этого потребовалось бы отдельно калибровать дискретные
автоморфизмы или ввести второй калибровочный кадровый множитель.
\end{nogocriterion}

\section{Решение развилки}

Получен не выбор по вкусу, а строгая трёхчастная граница.
\begin{enumerate}
 \item Составной третий момент является правильным тетраэдрическим
 параметром порядка и не требует ненулевого среднего кадра.
 \item Текущий конечный операторный носитель не содержит динамической
 семёрки спина три, поэтому этот момент пока остаётся эффективным.
 \item Упорядоченный четырёхосный кадр принадлежит текущей архитектуре,
 но полностью нарушает один калибровочный $SO(3)$ и оставляет только
 диагональную глобальную или граничную $A_4$.
\end{enumerate}

Если сохранить уже используемые $A_4$- и $\mathbb Z_3$-голономии как
настоящие дискретные калибровочные данные, минимальным продолжением
является не вторая семейная связность, а один канонический носитель
\begin{equation}
 \operatorname{Sym}^3_0(V_1)=V_3.
 \label{eq:v6-selected-minimal-spin-three-extension}
\end{equation}
Его представительный тип уникален, а потенциальный квадрат и нормировка
уже проверены предыдущей главой. Следующий гейт должен вывести этот
носитель как минимальное расширение конечной цепи и проверить, не приносит
ли он лишние поля, веса или аномалии.

\section{Вердикт}

\begin{protocol}
\begin{enumerate}
 \item чистый тетраэдрический третий момент: \textbf{получен};
 \item первый момент: \textbf{ноль};
 \item бесследовый второй момент: \textbf{ноль};
 \item стабилизатор третьего момента: \textbf{$A_4$};
 \item масса трёх семейных направлений: \textbf{положительна};
 \item спин три в $\operatorname{End}(V_1)$: \textbf{отсутствует};
 \item спин три в $\operatorname{Hom}(V_4,V_1)$: \textbf{отсутствует};
 \item спин три в $\operatorname{End}(K_{\rm fam})$: \textbf{отсутствует};
 \item спин три во внешнем $\operatorname{Sym}^3(V_1)$: \textbf{одна копия};
 \item правое $S_4$ как выведенная локальная калибровка:
 \textbf{нет};
 \item остаточная калибровочная группа упорядоченного кадра:
 \textbf{тривиальна};
 \item минимальная ветвь с сохранением $A_4\to\mathbb Z_3$:
 \textbf{носитель $V_3$};
 \item этот носитель уже встроен в конечного родителя: \textbf{нет};
 \item рождение материи: \textbf{не закрыто};
 \item следующий гейт: \textbf{минимальное конечное встраивание
 носителя спина три}.
\end{enumerate}
\end{protocol}

Машинный сертификат сохранён в
\path{s2t_v6_bosonic_defect_tetrahedral_gauge_frame_branch_decision_gate_results.json}.